Let
\[
 b=\sum_{a\in P_c}c_a=-\sum_{a\in N_c}c_a=\frac{L_c}{2}.
\]
The first equality uses \(\sum_a c_a=0\); the last uses the definition of \(L_c\).  Since \(c\neq0\), both \(P_c\) and \(N_c\) are nonempty and \(b>0\).  Given a two-arm binary schedule \(y=((Y_i(0),Y_i(1)))_{i=1}^n\), define a \(K\)-arm schedule \(\iota_c(y)\) by
\[
 t_{i,a}=\begin{cases}
 Y_i(1),&a\in P_c,\\
 Y_i(0),&a\in N_c,\\
 0,&a\notin S_c.
 \end{cases}
\]
This schedule is an element of \(\mathcal T_K^n\), including at the all-zero and all-one boundaries.  For every unit \(i\),
\[
 \sum_{a\in\mathcal A_K}c_at_{i,a}
 =Y_i(1)\sum_{a\in P_c}c_a+Y_i(0)\sum_{a\in N_c}c_a
 =b\{Y_i(1)-Y_i(0)\}.
\]
After summing over units and dividing by \(n\), this proves the displayed target identity.

Fix an arbitrary \(K\)-arm procedure \((\mathcal D,T)\), where \(T\) takes values in \([-b,b]\).  Define the deterministic sign-collapse map \(\beta:\mathcal A_K^n\to\{0,1\}^n\) coordinatewise by
\[
 \beta(A)_i=\mathbf 1\{A_i\in P_c\}.
\]
The pushforward
\[
 \mathcal D_2(B)=\sum_{A:\,\beta(A)=B}\mathcal D(A)
\]
is a legitimate two-arm assignment law: it is selected before the fixed binary schedule is revealed and may have arbitrary dependence across units.  Under assignment \(B\), let \(W_i=Y_i(B_i)\) denote the observed binary outcome.  For every \(A\) in the fiber \(\{A:\beta(A)=B\}\), form the synthetic \(K\)-arm observed outcome
\[
 \widetilde W_i(A,B,W)=
 \begin{cases}
 W_i,&A_i\in S_c,\\
 0,&A_i\notin S_c.
 \end{cases}
\]
If \(A_i\in P_c\), then \(B_i=1\) and \(W_i=Y_i(1)=t_{i,A_i}\).  If \(A_i\in N_c\), then \(B_i=0\) and \(W_i=Y_i(0)=t_{i,A_i}\).  If \(A_i\notin S_c\), the embedded potential outcome is the fixed value zero.  Thus \((A,\widetilde W)\) has exactly the data law of the original \(K\)-arm procedure on \(\iota_c(y)\).  In particular, the actual two-arm outcome observed at a unit whose latent arm is inactive is discarded; it is never supplied to \(T\).

For every \(B\) with \(\mathcal D_2(B)>0\), define
\[
 \eta(B,W)=
 \sum_{A:\,\beta(A)=B}
 \frac{\mathcal D(A)}{\mathcal D_2(B)}\,
 \frac{T(A,\widetilde W(A,B,W))}{b},
\]
and set \(\eta(B,W)=0\) when \(\mathcal D_2(B)=0\).  The weights depend only on the known design and \(B\), not on the response schedule.  Hence \(\eta\) is a deterministic function of legitimate two-arm observed data.  Moreover, \(T/b\in[-1,1]\), so convexity gives \(\eta\in[-1,1]\), the two-arm action space.

Let
\[
 \tau_2(y)=\frac1n\sum_{i=1}^n\{Y_i(1)-Y_i(0)\}.
\]
Conditional Jensen's inequality on each fiber of \(\beta\), followed by the target identity \(\tau_c(\iota_c(y))=b\tau_2(y)\), yields
\[
 \begin{aligned}
 R_{n,2}(\mathcal D_2,\eta;y)
 &=\mathbb E_{\mathcal D_2}\{\eta(B,W)-\tau_2(y)\}^2\\
 &\leq \mathbb E_{\mathcal D}
 \left[\left\{\frac{T(A,\widetilde W)}b-\tau_2(y)\right\}^2\right]\\
 &=\frac1{b^2}R_n(\mathcal D,T;\iota_c(y)).
 \end{aligned}
\]
Taking the maximum over all binary schedules on the left and then bounding the embedded schedules by all \(z\in\mathcal T_K^n\) gives
\[
 \max_yR_{n,2}(\mathcal D_2,\eta;y)
 \leq \frac1{b^2}\max_{z\in\mathcal T_K^n}R_n(\mathcal D,T;z).
\]
The left side is at least \(\rho_{n,2}\).  Since the original \((\mathcal D,T)\) was arbitrary,
\[
 \rho_n(c)\geq b^2\rho_{n,2}=C_0(c)\rho_{n,2}.
\]
No positivity condition is required: zero-probability fibers are assigned arbitrary estimator values, and the divisions above occur only when \(\mathcal D_2(B)>0\).  No rank or inverse condition is used.

Suppose next that \(|S_c|=2\).  Let \(a_+\) and \(a_-\) be the unique positive and negative active arms.  Zero-sum implies \(c_{a_+}=b\) and \(c_{a_-}=-b\).  Given any two-arm procedure \((\mathcal D_2,\eta)\), construct a \(K\)-arm procedure by replacing binary assignment one with \(a_+\), binary assignment zero with \(a_-\), never assigning an inactive arm, and reporting \(T=b\eta\).  For an arbitrary \(K\)-arm schedule, use \(Y_i(1)=t_{i,a_+}\) and \(Y_i(0)=t_{i,a_-}\).  Then
\[
 \tau_c(z)=b\tau_2(y),
 \qquad
 R_n(\mathcal D,T;z)=b^2R_{n,2}(\mathcal D_2,\eta;y).
\]
Inactive potential outcomes are neither assigned nor used.  Taking maxima and then the infimum over two-arm procedures proves \(\rho_n(c)\leq b^2\rho_{n,2}\), which combines with the converse to give equality.

It remains to establish the nonnegative improvement and its ceiling.  Under the contrast-weighted independent design, write
\[
 Z_i=\sum_{a\in S_c}\frac{c_a\mathbf 1\{A_i=a\}}{q_a^\star}
 \bigl(Y_i^{\mathrm{obs}}-1/2\bigr),
 \qquad \widetilde\tau=\frac1n\sum_{i=1}^nZ_i.
\]
On \(S_c\), \(q_a^\star=|c_a|/L_c>0\).  Independence across units, \(\sum_ac_a=0\), and binary outcomes give
\[
 \mathbb E Z_i=\sum_ac_at_{i,a},
 \qquad
 \mathbb E Z_i^2=\frac14\sum_{a\in S_c}\frac{c_a^2}{q_a^\star}
 =\frac{L_c^2}{4}=C_0(c).
\]
Consequently,
\[
 \mathbb E\{\widetilde\tau-\tau_c(z)\}^2
 =\frac1{n^2}\sum_{i=1}^n
 \left[C_0(c)-\left\{\sum_ac_at_{i,a}\right\}^2\right]
 \leq\frac{C_0(c)}n.
\]
The positive and negative coefficients each have absolute mass \(b=L_c/2\), so \(\tau_c(z)\in[-b,b]\).  Projection onto this interval cannot increase squared error, proving \(\rho_n(c)\leq C_0(c)/n\) and hence \(d_n(c)\geq0\).  Combining this upper bound with the embedded converse gives
\[
 d_n(c)=\frac{C_0(c)}n-\rho_n(c)
 \leq C_0(c)\{n^{-1}-\rho_{n,2}\}.
\]
The cited two-arm theorem states
\[
 \rho_{n,2}=n^{-1}-C_A n^{-4/3}+o(n^{-4/3}),
\]
which proves the asserted ceiling.  When \(|S_c|=2\), the exact equality \(\rho_n(c)=C_0(c)\rho_{n,2}\) makes the ceiling an equality.